\chapter{Congenital adrenal hyperplasia}
\index{Congenital adrenal hyperplasia}

\medskip
\noindent\textit{Educational information; seek specialist advice for individual care.}

\section*{What it is}
Congenital adrenal hyperplasia (CAH) is a group of inherited conditions in which the adrenal glands make too little of one or more hormones. The commonest form affects cortisol production and may also affect aldosterone and androgen levels. Severity ranges from forms presenting in infancy to milder forms diagnosed later.

\section*{Symptoms and diagnosis}
Severe forms may cause poor feeding, vomiting, dehydration, low blood pressure, low blood sugar, or unusual genital development in a newborn. Later features can include early puberty, rapid growth with advanced bone age, excess body hair, irregular periods, acne, or fertility difficulties. Diagnosis may involve hormone tests, electrolytes, genetic testing in selected situations, and assessment by an endocrinology team.

\section*{Treatment}
Treatment is individualized. It may include glucocorticoid replacement and, in salt-wasting forms, mineralocorticoid and salt replacement. Children need monitoring of growth, blood pressure, electrolytes, hormones, and treatment effects. Decisions about genital surgery, fertility, mental health, and genetic counselling should be discussed carefully with the person and family.

\section*{Preventing adrenal crisis}
People with CAH and their caregivers should have a written sick-day plan, an emergency identification card or bracelet, and training in emergency hydrocortisone when prescribed. Illness, injury, surgery, or inability to keep medicines down may require an increased glucocorticoid dose or injection. Follow the individual plan; do not change doses without clinical advice.

\section*{When to Seek Medical Care}
Get emergency help for repeated vomiting, severe weakness, confusion, fainting, dehydration, low blood pressure, or inability to take prescribed steroid medicine. If an emergency injection has been prescribed and the person cannot keep oral medicine down, use it as trained and seek emergency care.

\section*{Sources}
\begin{itemize}
\item Endocrine Society patient information: \url{https://www.endocrine.org/patient-engagement/endocrine-library/congenital-adrenal-hyperplasia}
\item Endocrine Society clinical practice guideline resources: \url{https://www.endocrine.org/clinical-practice-guidelines/congenital-adrenal-hyperplasia-guideline-resources}
\end{itemize}
